\documentclass[onecolumn,12pt]{IEEEtran}
\usepackage[utf8]{inputenc}
\usepackage[T1]{fontenc}
\usepackage[spanish,es-tabla]{babel}
\usepackage{booktabs}
\usepackage{longtable}
\usepackage{array}
\usepackage{graphicx}
\usepackage{amsmath}
\usepackage{tikz}
\usepackage[hidelinks]{hyperref}
\AtBeginDocument{\spanishplainpercent}
\renewcommand{\IEEEkeywordsname}{Palabras clave}

\title{FORM-SIM-01: Simulación de priorización y negociación de requisitos\\
para matrícula y calificaciones en una unidad educativa}
\author{\IEEEauthorblockN{Mendoza Palma Allan Jeremy, Rochina Rochina Freddy David, Sánchez Centeno Roselyn Andreina}
\IEEEauthorblockA{Facultad de Ciencias de la Computación, Carrera de Ingeniería de Software\\
Universidad Técnica Estatal de Quevedo\\
Ingeniería de Requerimientos (ISR-401)}}

\begin{document}
\maketitle
\begin{center}
\small
\begin{tabular}{@{}p{0.27\textwidth}p{0.59\textwidth}@{}}
\toprule
\textbf{Dato informativo} & \textbf{Descripción}\\
\midrule
Actividad & FORM-SIM-01, simulación de priorización y negociación entre roles\\
Modalidad & Formativa, trabajo grupal\\
Asignatura y paralelo & Ingeniería de Requerimientos (ISR-401), cuarto nivel, paralelo B\\
Docente & Gleiston Cicerón Guerrero Ulloa, PhD\\
Unidad & Unidad 2: Elicitación y análisis de requisitos\\
Período & SPA 2026--2027\\
Dominio & Unidad educativa: matrícula y calificaciones\\
Cierre & Sábado 28 de noviembre de 2026, 23:59\\
Repositorio & \url{https://github.com/AllanMendoza09/IR-trabajos-grupales-MRS}\\
\bottomrule
\end{tabular}
\end{center}
\vspace{0.8em}

\begin{abstract}
Se establece una estructura reproducible para analizar y negociar la prioridad de requisitos de un sistema de matrícula y calificaciones. Se definen criterios de especificación, el método MoSCoW, un instrumento Kano de ocho características, una votación de cien puntos por participante y procedimientos de comparación de rangos. El análisis empírico se desarrollará con respuestas obtenidas de participantes adultos, identificados mediante códigos, sin publicar datos personales.
\end{abstract}

\begin{IEEEkeywords}
Ingeniería de requisitos, MoSCoW, Kano, votación ponderada, Spearman, educación.
\end{IEEEkeywords}

\noindent\textbf{Propósito:} aplicar técnicas de priorización y negociación entre roles a un conjunto de requisitos de matrícula y calificaciones, documentando criterios y resultados reproducibles.

\IfFileExists{secciones/00_fundamento.tex}{\section{Fundamento y delimitación metodológica}
La Ingeniería de Requisitos transforma necesidades de los involucrados en declaraciones evaluables y negociables. La especificación individual requiere condiciones explícitas y criterios que permitan determinar su aceptación, mientras que la priorización selecciona qué capacidades aportan valor dentro de restricciones de esfuerzo y alcance \cite{iso29148,pohl2025,wiegers2013}.

\subsection{Sistema de estudio y unidad de análisis}
El dominio es una unidad educativa con procesos de matrícula, asignación de estudiantes a cursos, registro de calificaciones, consulta del historial y supervisión académica. La unidad de priorización es cada requisito identificado con un código único. Se diferencian requisitos funcionales (RF) de requisitos de calidad (RNF). El conjunto de trabajo comprende veinte requisitos: doce RF y ocho RNF con criterios cuantificados. Los umbrales de los RNF representan objetivos propuestos de diseño y no resultados de mediciones de una institución concreta.

\subsection{Técnicas de priorización y actores}
\textbf{MoSCoW} clasifica cada requisito como Must, Should, Could o Won't en la iteración considerada. La categoría Must se contrasta con el esfuerzo total estimado por el equipo. Su límite es
\begin{equation}
\frac{\sum_{i\in M}e_i}{\sum_{i=1}^{20}e_i}\leq 0{,}60,\qquad e_i>0.
\label{eq:must}
\end{equation}
La clasificación se justifica mediante una frase verificable sobre el valor, la dependencia o el riesgo que resuelve cada requisito; no sustituye el análisis de factibilidad.

\textbf{Kano} contrasta la reacción del participante ante la presencia y la ausencia de cada una de ocho características. Para cada característica se registran dos respuestas en una escala uniforme de cinco opciones. La matriz de evaluación produce clases A (atractiva), O (unidimensional), M (básica), I (indiferente), R (inversa) y Q (contradictoria). Las ocho características seleccionadas forman un subconjunto del catálogo; por tanto, Kano no origina por sí solo un orden total para los veinte requisitos.

\textbf{Votación ponderada} permite que cada una de tres personas distribuya exactamente cien puntos sobre los veinte requisitos. Se conserva cada distribución individual para revisar la legitimidad de la suma de puntos; no se imputan puntos ni preferencias a participantes no consultados.

\subsection{Rangos, empates y correlación}
Los rangos de MoSCoW agrupan requisitos por categoría y asignan el promedio de las posiciones ocupadas en caso de empate. La votación ordena los requisitos por puntos totales y utiliza el mismo procedimiento de rangos medios. En lugar de la fórmula abreviada que presupone ausencia de empates, se calcula Spearman como correlación de Pearson entre los vectores de rangos \(r_i\) y \(s_i\):
\begin{equation}
\rho=\frac{\sum_i(r_i-\bar r)(s_i-\bar s)}{\sqrt{\sum_i(r_i-\bar r)^2\sum_i(s_i-\bar s)^2}}.
\label{eq:spearman}
\end{equation}
Si alguno de los órdenes carece de variación, no se informa coeficiente porque el denominador sería cero. La coincidencia estadística no sustituye la deliberación, ni permite afirmar causalidad.

\subsection{Protección y reproducibilidad}
Se prevé la participación de cinco adultos en Kano y tres adultos distintos en votación. Sus respuestas se vincularán exclusivamente a los códigos P01--P08. Las evidencias originales y consentimientos se conservarán fuera del repositorio público; el manifiesto documentará únicamente códigos, técnicas, fechas, nombres de evidencia privada y huellas SHA-256. La publicación se limitará a resultados anonimizados y artefactos reproducibles, en concordancia con la protección de datos y la ética profesional \cite{lopdp2021,acm2018}.
}{}
\IfFileExists{secciones/01_requisitos.tex}{\section{Catálogo de requisitos candidatos}
\label{sec:catalogo-requisitos}

El caso de estudio corresponde a los procesos de matrícula y registro de calificaciones de una unidad educativa. Se contemplan estudiantes, representantes, docentes, personal de secretaría, coordinación académica y autoridades como partes interesadas en los procesos escolares. Los requisitos se especifican como comportamientos y condiciones observables mediante el verbo modal \emph{deberá}, de conformidad con el enfoque de especificación individual \cite{iso29148,wiegers2013}.

Se analizan doce requisitos funcionales (RF) y ocho requisitos de calidad (RNF). Las cifras utilizadas en los RNF son \emph{umbrales propuestos}, no mediciones de una unidad educativa específica. El alcance está limitado a las funciones académicas de matrícula y calificaciones.

\begingroup
\footnotesize
\renewcommand{\arraystretch}{1.12}
\begin{longtable}{@{}p{0.12\linewidth}p{0.83\linewidth}@{}}
\caption{Catálogo de veinte requisitos candidatos para la unidad educativa}\label{tab:catalogo-veinte}\\
\toprule
\textbf{ID} & \textbf{Requisito propuesto} \\
\midrule
\endfirsthead
\multicolumn{2}{@{}l}{\textit{Continuación del catálogo de requisitos}}\\
\toprule
\textbf{ID} & \textbf{Requisito propuesto} \\
\midrule
\endhead
\bottomrule
\endlastfoot
\textbf{RF-01} & El sistema deberá registrar un estudiante con identificador institucional único cuando la secretaría confirme los datos de matrícula y la identificación no exista, con el fin de crear su expediente académico. \\
\textbf{RF-02} & El sistema deberá autenticar a un usuario mediante credenciales institucionales cuando solicite ingresar a la plataforma, con el fin de habilitar las funciones correspondientes a su perfil. \\
\textbf{RF-03} & El sistema deberá mostrar los paralelos ofertados en el período lectivo seleccionado cuando un usuario autorizado consulte la oferta académica, con el fin de identificar opciones de matrícula. \\
\textbf{RF-04} & El sistema deberá registrar la matrícula de un estudiante en el paralelo seleccionado cuando exista un cupo disponible y el estudiante cumpla las condiciones académicas definidas, con el fin de formalizar su inscripción. \\
\textbf{RF-05} & El sistema deberá rechazar una solicitud de matrícula en un paralelo cuando el número de matriculados alcance su capacidad definida, con el fin de evitar sobreasignación de cupos. \\
\textbf{RF-06} & El sistema deberá generar un comprobante con código, estudiante, paralelo y período lectivo cuando la matrícula quede registrada, con el fin de acreditar la inscripción. \\
\textbf{RF-07} & El sistema deberá asociar un docente autorizado a una asignatura de un paralelo cuando la coordinación confirme su distribución académica, con el fin de habilitar el registro de notas. \\
\textbf{RF-08} & El sistema deberá registrar la calificación numérica de una evaluación para un estudiante matriculado cuando el docente asignado confirme el valor dentro de la escala institucional configurada, con el fin de mantener el registro académico. \\
\textbf{RF-09} & El sistema deberá registrar una nueva versión de una calificación con la justificación de modificación cuando el responsable autorizado apruebe su corrección, con el fin de conservar el historial de cambios. \\
\textbf{RF-10} & El sistema deberá mostrar al estudiante autenticado sus calificaciones publicadas del período seleccionado cuando consulte su historial, con el fin de conocer su desempeño académico. \\
\textbf{RF-11} & El sistema deberá generar el boletín de un estudiante con asignaturas y calificaciones publicadas cuando secretaría solicite el documento para un período lectivo, con el fin de respaldar su trayectoria académica. \\
\textbf{RF-12} & El sistema deberá generar un reporte de matriculados por paralelo y período cuando la dirección académica seleccione un período lectivo, con el fin de supervisar la ocupación de cupos. \\
\textbf{RNF-01} & El sistema deberá completar al menos el 95 \% de los registros de matrícula en un máximo de 3 segundos cuando existan 100 solicitudes simultáneas sobre una base de 3000 estudiantes, con el fin de mantener tiempos controlados en el inicio de período. \\
\textbf{RNF-02} & El sistema deberá mantener una disponibilidad mensual mínima del 99,5 \% del portal cuando se evalúe un mes calendario de operación, incluyendo mantenimientos, con el fin de asegurar acceso regular a servicios académicos. \\
\textbf{RNF-03} & El sistema deberá impedir el acceso a calificaciones individuales de un estudiante desde cuentas no autorizadas en el 100 \% de los casos cuando se ejecute una batería de 50 solicitudes negativas con perfiles de estudiante, docente y secretaría, con el fin de proteger información académica personal. \\
\textbf{RNF-04} & El sistema deberá permitir que al menos el 90 \% de una muestra de 20 usuarios adultos complete una matrícula simulada en no más de 5 minutos cuando utilicen por primera vez la interfaz sin ayuda, con el fin de favorecer la autogestión. \\
\textbf{RNF-05} & El sistema deberá mostrar sin desplazamiento horizontal las pantallas de consulta de matrícula y notas en al menos cinco resoluciones de pantalla entre 360 y 768 píxeles cuando se pruebe en navegadores modernos, con el fin de facilitar su uso en teléfonos. \\
\textbf{RNF-06} & El sistema deberá generar al menos una copia íntegra de las bases de matrícula y calificaciones cada 24 horas cuando se encuentre en operación, con el fin de limitar la pérdida de datos ante fallos. \\
\textbf{RNF-07} & El sistema deberá restaurar el servicio de consulta de calificaciones en un máximo de 60 minutos cuando se simule una caída completa del servidor en un ambiente de recuperación, con el fin de reducir la interrupción académica. \\
\textbf{RNF-08} & El sistema deberá registrar el 100 \% de las modificaciones de calificaciones con identificador de actor, momento, valor anterior y valor nuevo cuando se ejecuten 50 cambios autorizados de prueba, con el fin de permitir auditorías internas. \\
\end{longtable}
\endgroup

\subsection{Criterios de revisión de los requisitos}

La revisión de cada requisito considera el comportamiento solicitado, su condición de activación, el propósito, la autorización del actor y la posibilidad de comprobarlo mediante pruebas. En los requisitos de calidad se debe verificar explícitamente la métrica, el umbral, la unidad y el escenario de medición. El archivo \texttt{datos/requisitos.csv} contiene el código, tipo, nombre, redacción y forma propuesta de verificación, lo cual permite cotejar este catálogo con los datos empleados en las técnicas de priorización.

Las políticas de matrícula, promoción, escalas de calificaciones, cupos y roles de acceso deberán contrastarse con la normativa y la organización de la unidad educativa donde se aplique el ejercicio. Esta necesidad de validación no implica que se hayan efectuado consultas institucionales.
}{}
\IfFileExists{secciones/02_moscow.tex}{\section{Priorización con MoSCoW}
\label{sec:moscow}

\subsection{Método}
\label{subsec:moscow-metodo}

La técnica MoSCoW asigna a cada requisito una de cuatro categorías que expresan un acuerdo de prioridad y de alcance. Las categorías no constituyen una medición matemática de la importancia de un requisito. Se definen así:

\begin{itemize}
  \item \textbf{Must:} requisito indispensable para que la solución cumpla su propósito mínimo acordado.
  \item \textbf{Should:} requisito de alta importancia que puede aplazarse sin impedir el funcionamiento mínimo acordado.
  \item \textbf{Could:} requisito deseable que aporta valor, pero tiene menor prioridad relativa.
  \item \textbf{Won't:} requisito que se excluye del alcance de la entrega actual, aunque puede reconsiderarse posteriormente.
\end{itemize}

Para asignar la categoría se consideraron la necesidad que atiende el requisito, su contribución al propósito mínimo del sistema, sus dependencias con otros requisitos y las consecuencias de no implementarlo. Cada categoría cuenta con una justificación de una línea.

\subsection{Estimación del esfuerzo}
\label{subsec:moscow-esfuerzo}

El esfuerzo se expresa en puntos relativos de esfuerzo, con la escala 1, 2, 3, 5, 8 y 13. Cada requisito se compara con los demás y recibe un valor que representa su esfuerzo previsto en relación con el resto. Los valores son estimaciones del equipo y no corresponden a mediciones históricas ni a duraciones oficiales. La misma unidad se usa en los veinte requisitos y en todos los cálculos.

\subsection{Tabla de priorización}
\label{subsec:moscow-tabla}

La tabla~\ref{tab:moscow} presenta la categoría, el esfuerzo y la justificación de cada requisito.

{\small
\setlength{\tabcolsep}{3pt}
\begin{longtable}{@{}>{\raggedright\arraybackslash}p{0.10\linewidth}>{\raggedright\arraybackslash}p{0.10\linewidth}>{\raggedright\arraybackslash}p{0.12\linewidth}>{\raggedright\arraybackslash}p{0.62\linewidth}@{}}
\caption{Priorización MoSCoW de los requisitos.}\label{tab:moscow}\\
\hline
\textbf{ID} & \textbf{Categoría} & \textbf{Esfuerzo (puntos)} & \textbf{Justificación}\\
\hline
\endfirsthead
\caption[]{Priorización MoSCoW de los requisitos (continuación).}\\
\hline
\textbf{ID} & \textbf{Categoría} & \textbf{Esfuerzo (puntos)} & \textbf{Justificación}\\
\hline
\endhead
\hline
\endfoot
RF-01 & Must & 5 & Sin el expediente del estudiante no existe base para matricular ni calificar\\
RF-02 & Must & 5 & Sin autenticación no se pueden habilitar perfiles ni proteger los datos\\
RF-03 & Must & 3 & El estudiante necesita ver los paralelos ofertados para poder matricularse\\
RF-04 & Must & 8 & Es la función central del proceso de matrícula\\
RF-05 & Must & 5 & Evita matrículas por encima de la capacidad del paralelo\\
RF-06 & Should & 3 & La inscripción puede acreditarse de otra forma mientras se implementa el comprobante\\
RF-07 & Must & 5 & Sin docente asignado a la asignatura no se pueden registrar notas\\
RF-08 & Must & 5 & Es el núcleo del registro académico de calificaciones\\
RF-09 & Should & 5 & Las correcciones pueden gestionarse fuera del sistema sin impedir el registro básico\\
RF-10 & Must & 3 & Es el servicio central de consulta para el estudiante y depende de RF-08\\
RF-11 & Should & 5 & El boletín puede elaborarse manualmente sin impedir el funcionamiento mínimo\\
RF-12 & Could & 3 & Aporta supervisión de cupos, pero no condiciona la matrícula ni las notas\\
RNF-01 & Should & 5 & Mejora el inicio de período; el mínimo funciona con tiempos menos estrictos\\
RNF-02 & Should & 8 & Valioso, pero de alto costo y no define el mínimo viable\\
RNF-03 & Must & 5 & Protege la información académica personal frente a accesos no autorizados\\
RNF-04 & Could & 3 & Deseable; su ausencia no bloquea el servicio\\
RNF-05 & Could & 5 & Deseable; el acceso desde computador cubre el mínimo viable\\
RNF-06 & Should & 3 & Limita la pérdida de datos; no define el mínimo viable\\
RNF-07 & Could & 5 & Reduce la interrupción tras una caída; puede abordarse en una etapa posterior\\
RNF-08 & Should & 3 & Respalda las correcciones (RF-09), que también son Should\\
\end{longtable}
}

\subsection{Verificación del límite de esfuerzo del grupo Must}
\label{subsec:moscow-limite}

El esfuerzo estimado del grupo Must no debe superar el 60\,\% del esfuerzo estimado total. Se verifica con las expresiones

\begin{equation}
E_{\mathrm{total}}=\sum_{i=1}^{20}E_i,\qquad
E_{\mathrm{Must}}=\sum_{i\in\mathrm{Must}}E_i,\qquad
P_{\mathrm{Must}}=\frac{E_{\mathrm{Must}}}{E_{\mathrm{total}}}\times 100,
\label{eq:moscow-limite}
\end{equation}

donde $E_i$ es el esfuerzo estimado del requisito $i$, $E_{\mathrm{total}}$ la suma del esfuerzo de los veinte requisitos, $E_{\mathrm{Must}}$ la suma del esfuerzo de los requisitos Must y $P_{\mathrm{Must}}$ el porcentaje del esfuerzo que corresponde a Must. El límite se cumple cuando $P_{\mathrm{Must}}\leq 60\,\%$.

Con los valores de la tabla~\ref{tab:moscow}, $E_{\mathrm{total}}=92$ puntos y $E_{\mathrm{Must}}=44$ puntos, de modo que $P_{\mathrm{Must}}=44/92\times 100\approx 47,83\,\%$. El valor es inferior al 60\,\% y el límite se cumple. El esfuerzo de Should suma 32 puntos, el de Could 16 puntos y el de Won't 0 puntos.
}{}
\IfFileExists{secciones/03_kano.tex}{\section{Protocolo de simulación y registro de evidencias}
\label{sec:protocolo-simulacion}

La simulación permite contrastar las preferencias de los interesados ante funcionalidades del sistema de matrícula y calificaciones de una unidad educativa. El instrumento Kano establece una pregunta funcional (presencia de la característica) y una disfuncional (ausencia), usando cinco opciones comunes \cite{pohl2025,wiegers2013}. El protocolo describe una aplicación prevista; no constituye un informe de encuestas ya ejecutadas.

\subsection{Características de evaluación Kano}

Se seleccionaron ocho características relacionadas con los veinte requisitos candidatos. La correspondencia entre el código Kano y el código de requisito se conserva en \path{datos/caracteristicas_kano.csv}.

\begingroup
\small
\begin{longtable}{@{}p{0.15\linewidth}p{0.18\linewidth}p{0.61\linewidth}@{}}
\caption{Correspondencia entre características Kano y requisitos escolares}\label{tab:kano-caracteristicas}\\
\toprule
\textbf{Código} & \textbf{Requisito} & \textbf{Característica} \\
\midrule
K01 & RF-04 & Matrícula en línea \\
K02 & RF-06 & Comprobante digital de matrícula \\
K03 & RF-08 & Registro de calificaciones \\
K04 & RF-10 & Consulta personal de calificaciones \\
K05 & RF-11 & Boletín académico digital \\
K06 & RNF-01 & Rapidez de matrícula \\
K07 & RNF-04 & Facilidad de uso \\
K08 & RNF-05 & Compatibilidad con móvil \\
\bottomrule
\end{longtable}
\endgroup

\subsection{Procedimiento de aplicación}

Se prevé aplicar el instrumento a cinco personas adultas distintas, vinculadas con actividades de matrícula o seguimiento escolar, identificadas mediante los códigos P01 a P05. Los roles pertinentes al dominio incluyen secretaría, docentes, representantes, autoridades y estudiantes adultos. Para cada característica, el participante responderá la versión funcional y la disfuncional de una pregunta. El instrumento incluye dieciséis preguntas y se encuentra en \path{datos/instrumento_kano.csv}.

Las respuestas se registrarán en el archivo CSV de Kano. Cada fila identificará el código del participante, la característica evaluada y sus dos respuestas. La estructura de los campos está descrita en \path{datos/ESQUEMA.md}. Los cinco participantes deben aportar ocho pares válidos cada uno, para un mínimo de cuarenta registros. El análisis posterior utilizará la matriz de Kano y distinguirá las categorías básica, unidimensional, atractiva, indiferente, inversa y contradictoria. No se imputarán respuestas ausentes.

\subsection{Votación ponderada y negociación}

Se prevé una segunda técnica con tres personas adultas diferentes de las participantes en Kano, identificadas como P06, P07 y P08. Cada participante distribuirá exactamente cien puntos entre los veinte requisitos del catálogo; sus votos se registrarán en \path{datos/votos.csv}. La clasificación MoSCoW será discutida por el equipo y documentada en \path{datos/moscow.csv}, incluyendo para cada requisito una prioridad, una estimación positiva de esfuerzo y una justificación. El esfuerzo acumulado del grupo Must no podrá superar el sesenta por ciento del esfuerzo total acordado. La comparación y las decisiones finales no deben atribuirse a participantes sin datos recopilados.

\subsection{Consentimiento, custodia y manifiesto}

La invitación a participar debe explicar el propósito académico, el carácter voluntario del ejercicio, el tratamiento de respuestas anonimizadas y la custodia privada de los respaldos. No se publicarán formularios firmados, audios, videos, nombres, teléfonos, correos personales ni documentos de identificación. El equipo conservará los originales en una carpeta privada accesible a las personas autorizadas y registrará sus metadatos en \path{evidencias/MANIFIESTO.md}, de acuerdo con prácticas prudentes de protección de información personal \cite{lopdp2021,acm2018}.

Cada fila del manifiesto se incorporará únicamente al disponer del archivo original. Registrará código del participante, rol en el dominio educativo, técnica, fecha de recolección, nombre del archivo privado y huella SHA-256 obtenida de sus bytes. El manifiesto público no contiene entradas de participantes hasta que existan evidencias efectivamente obtenidas.
}{}
\IfFileExists{secciones/04_votacion.tex}{\section{Votación ponderada: registro y resultados}
\label{sec:votacion}

\subsection{Datos y control de calidad}

El 27 de noviembre de 2026 se registraron tres formularios identificados como P06, P07 y P08, distintos de los cinco códigos de Kano. Cada formulario contiene una asignación de puntos enteros no negativos para los veinte requisitos (RF-01 a RF-12 y RNF-01 a RNF-08). Las tres sumas individuales son exactamente cien puntos, por lo que el conjunto agregado totaliza trescientos puntos.

Para obtener el puntaje total de un requisito se suman sus tres asignaciones. El rango 1 corresponde al mayor total. Los empates reciben el rango medio de las posiciones que ocupan; así, dos requisitos empatados en las posiciones 6 y 7 tienen rango 6,5. Se conservan los valores individuales en \path{datos/votos.csv} y el cálculo verificable en \path{datos/calculos/orden_votacion.csv}.

\subsection{Orden obtenido}

\begingroup
\small
\begin{longtable}{@{}p{0.13\linewidth}p{0.53\linewidth}rr@{}}
\caption{Votación ponderada y rango medio}\label{tab:resultados-votacion}\\
\toprule
\textbf{ID} & \textbf{Requisito} & \textbf{Puntos} & \textbf{Rango}\\
\midrule
\endfirsthead
\toprule
\textbf{ID} & \textbf{Requisito} & \textbf{Puntos} & \textbf{Rango}\\
\midrule
\endhead
\bottomrule
\endlastfoot
RF-04 & Matrícula & 31 & 1,0\\
RF-08 & Registro de calificación & 25 & 2,0\\
RF-05 & Control de cupos & 24 & 3,0\\
RF-01 & Registro de estudiante & 21 & 4,0\\
RF-10 & Consulta de notas & 19 & 5,0\\
RF-07 & Asignación docente & 18 & 6,5\\
RF-12 & Reporte de matrícula & 18 & 6,5\\
RF-09 & Corrección de calificación & 17 & 8,0\\
RF-02 & Autenticación & 16 & 9,0\\
RNF-03 & Privacidad de calificaciones & 15 & 10,0\\
RF-06 & Comprobante de matrícula & 13 & 11,0\\
RF-11 & Boletín académico & 12 & 12,5\\
RNF-02 & Disponibilidad & 12 & 12,5\\
RF-03 & Oferta de cursos & 11 & 14,5\\
RNF-01 & Rendimiento de matrícula & 11 & 14,5\\
RNF-04 & Facilidad de uso de matrícula & 9 & 16,5\\
RNF-05 & Compatibilidad móvil & 9 & 16,5\\
RNF-06 & Copia de seguridad & 8 & 18,0\\
RNF-07 & Recuperación del servicio & 6 & 19,0\\
RNF-08 & Trazabilidad de modificaciones & 5 & 20,0\\
\end{longtable}
\endgroup


RF-04 (matrícula) obtuvo 31 puntos; RF-08 (registro de calificaciones), 25; y RF-05 (control de cupos), 24. El primer requisito de calidad en la ordenación fue RNF-03 (privacidad de calificaciones), con 15 puntos. La posición relativa deriva únicamente de estas tres papeletas, no de una valoración de toda la comunidad educativa.
}{}
\IfFileExists{secciones/05_comparacion.tex}{\section{Comparación de los tres procedimientos}
\label{sec:comparacion}

\subsection{MoSCoW y votación ponderada}

MoSCoW clasificó veinte requisitos en nueve Must, siete Should y cuatro Could; ninguno recibió categoría Won't. Los valores comunicados por el equipo corresponden a 44 puntos de esfuerzo Must, 32 Should y 16 Could, con un total de 92 puntos relativos. El grupo Must representa el 47,83\,\% y cumple la restricción de la ecuación~\eqref{eq:must}.

Se asignaron rangos medios a los requisitos empatados dentro de cada grupo MoSCoW: Must ocupa las posiciones 1 a 9 (rango 5), Should las posiciones 10 a 16 (rango 13) y Could las posiciones 17 a 20 (rango 18,5). La votación se ordenó por puntos descendentes, también con rangos medios cuando las puntuaciones coinciden. Con los veinte pares de rangos, la correlación de Spearman, calculada como correlación de Pearson entre rangos, fue \(\rho=0{,}6605\). Se trata de una asociación positiva apreciable, pero no de una equivalencia entre las técnicas. No se infiere causalidad ni se generaliza a una población escolar: la votación recoge tres participantes.

\begingroup
\small
\begin{longtable}{@{}p{0.14\linewidth}p{0.21\linewidth}rrr@{}}
\caption{Contraste entre MoSCoW y votación}\label{tab:comparacion-moscow-votos}\\
\toprule
\textbf{ID} & \textbf{MoSCoW} & \textbf{R. MoSCoW} & \textbf{Puntos} & \textbf{R. votos}\\
\midrule
\endfirsthead
\toprule
\textbf{ID} & \textbf{MoSCoW} & \textbf{R. MoSCoW} & \textbf{Puntos} & \textbf{R. votos}\\
\midrule
\endhead
\bottomrule
\endlastfoot
RF-04 & Must & 5,0 & 31 & 1,0\\
RF-08 & Must & 5,0 & 25 & 2,0\\
RF-05 & Must & 5,0 & 24 & 3,0\\
RF-01 & Must & 5,0 & 21 & 4,0\\
RF-10 & Must & 5,0 & 19 & 5,0\\
RF-07 & Must & 5,0 & 18 & 6,5\\
RF-12 & Could & 18,5 & 18 & 6,5\\
RF-09 & Should & 13,0 & 17 & 8,0\\
RF-02 & Must & 5,0 & 16 & 9,0\\
RNF-03 & Must & 5,0 & 15 & 10,0\\
RF-06 & Should & 13,0 & 13 & 11,0\\
RF-11 & Should & 13,0 & 12 & 12,5\\
RNF-02 & Should & 13,0 & 12 & 12,5\\
RF-03 & Must & 5,0 & 11 & 14,5\\
RNF-01 & Should & 13,0 & 11 & 14,5\\
RNF-04 & Could & 18,5 & 9 & 16,5\\
RNF-05 & Could & 18,5 & 9 & 16,5\\
RNF-06 & Should & 13,0 & 8 & 18,0\\
RNF-07 & Could & 18,5 & 6 & 19,0\\
RNF-08 & Should & 13,0 & 5 & 20,0\\
\end{longtable}
\endgroup


Entre las coincidencias, RF-04 (matrícula, 31 puntos), RF-08 (registro de calificaciones, 25) y RF-05 (control de cupos, 24) son Must y ocupan las tres primeras posiciones de votación. En cambio, RF-12 (reporte de matrícula) recibe 18 puntos y comparte el rango 6,5 en la votación, aunque está clasificado como Could. RF-03 es Must, pero obtiene 11 puntos y comparte el rango 14,5. Estas diferencias justifican una negociación del alcance guiada por dependencias y esfuerzo, no solamente por la suma de puntos.

\subsection{Kano frente a la votación}

Kano devuelve una categoría para ocho características, no un ranking de veinte requisitos. Se observaron dos características indiferentes (K01 y K08), cuatro unidimensionales (K02 a K05) y dos atractivas (K06 y K07). En la tabla siguiente se muestran solo esas ocho correspondencias, sin imputar una categoría a los doce requisitos no evaluados por Kano.

\begingroup
\small
\begin{longtable}{@{}p{0.14\linewidth}p{0.40\linewidth}rrc@{}}
\caption{Características Kano frente a la votación}\label{tab:comparacion-kano-votos}\\
\toprule
\textbf{ID} & \textbf{Característica} & \textbf{Puntos} & \textbf{Rango} & \textbf{Kano}\\
\midrule
\endfirsthead
\toprule
\textbf{ID} & \textbf{Característica} & \textbf{Puntos} & \textbf{Rango} & \textbf{Kano}\\
\midrule
\endhead
\bottomrule
\endlastfoot
RF-04 & Matrícula & 31 & 1,0 & I\\
RF-06 & Comprobante de matrícula & 13 & 11,0 & O\\
RF-08 & Registro de calificación & 25 & 2,0 & O\\
RF-10 & Consulta de notas & 19 & 5,0 & O\\
RF-11 & Boletín académico & 12 & 12,5 & O\\
RNF-01 & Rendimiento de matrícula & 11 & 14,5 & A\\
RNF-04 & Facilidad de uso de matrícula & 9 & 16,5 & A\\
RNF-05 & Compatibilidad móvil & 9 & 16,5 & I\\
\end{longtable}
\endgroup


RF-04 fue el requisito mejor puntuado de la votación (31 puntos), pero su característica K01 obtuvo tres de cinco respuestas indiferentes. RF-08 alcanzó 25 puntos y su característica K03 fue unidimensional. Esta diferencia metodológica no se interpreta como contradicción de las personas participantes, porque las técnicas preguntan por aspectos diferentes. El tamaño de las muestras (cinco formularios Kano y tres votantes) restringe la interpretación a los registros analizados. En la tabulación Kano publicada se repitió el código P04 y no aparece P03; la copia exportada proporcionada por el equipo sí registra P03. No se alteran los archivos de respuestas de otros integrantes en esta sección.
}{}
\IfFileExists{secciones/06_decision.tex}{\section{Propuesta de alcance y criterios de exclusión}
\label{sec:decision}

La propuesta de primera versión aplica una regla explícita de corte: incorporar los nueve requisitos clasificados como Must en el archivo MoSCoW del equipo, y ordenar su atención interna según los puntos obtenidos en la votación. El esfuerzo estimado de este conjunto es de 44 puntos sobre 92 (47,83\,\%), por debajo del máximo de 60\,\%. La propuesta recoge los datos disponibles y requiere la ratificación del equipo antes de registrarse como decisión definitiva.

\subsection{Requisitos propuestos para la primera versión}

\begingroup\small
\begin{longtable}{@{}p{0.10\linewidth}p{0.16\linewidth}p{0.52\linewidth}r@{}}
\caption{Orden propuesto de requisitos Must para la primera versión}\label{tab:propuesta-primera-version}\\
\toprule
\textbf{Orden} & \textbf{ID} & \textbf{Requisito} & \textbf{Votos}\\
\midrule\endfirsthead
\toprule\textbf{Orden}&\textbf{ID}&\textbf{Requisito}&\textbf{Votos}\\\midrule\endhead
\bottomrule\endlastfoot
1 & RF-04 & Matrícula & 31\\
2 & RF-08 & Registro de calificaciones & 25\\
3 & RF-05 & Control de cupos & 24\\
4 & RF-01 & Registro de estudiante & 21\\
5 & RF-10 & Consulta de notas & 19\\
6 & RF-07 & Asignación docente & 18\\
7 & RF-02 & Autenticación & 16\\
8 & RNF-03 & Privacidad de calificaciones & 15\\
9 & RF-03 & Oferta de cursos & 11\\
\end{longtable}
\endgroup

El orden de lectura no es necesariamente el orden técnico de implementación. Autenticación (RF-02) y registro de estudiantes (RF-01), por ejemplo, constituyen condiciones de operación de otros requisitos y deben coordinarse como dependencias antes de programar entregas.

\subsection{Requisitos fuera del alcance propuesto de la primera versión}

Los siete Should se aplazan bajo la regla de corte por no ser indispensables en la versión mínima acordada en MoSCoW: RF-06, RF-09, RF-11, RNF-01, RNF-02, RNF-06 y RNF-08. Suman 32 puntos de esfuerzo. Los cuatro Could quedan fuera de esta versión por su menor prioridad relativa: RF-12, RNF-04, RNF-05 y RNF-07, con 16 puntos de esfuerzo. Esta diferenciación mantiene los veinte requisitos rastreables y no equivale a descartar definitivamente once requisitos.

No se registró ningún Won't. Por tanto, no existen requisitos formalmente descartados mediante esa categoría en la clasificación disponible. Una decisión posterior de exclusión o incorporación deberá quedar justificada por ID, prioridad, esfuerzo y aprobación del equipo. La diferencia entre RF-12 (Could con 18 votos) y RF-03 (Must con 11 votos) merece discusión explícita en esa negociación.
}{}
\IfFileExists{secciones/07_uso_ia.tex}{\section{Declaración de uso de inteligencia artificial}
Durante la preparación técnica de esta primera estructura, Mendoza Palma Allan Jeremy solicitó asistencia mediante ChatGPT (modelo GPT-6) para organizar el proyecto, redactar versiones iniciales de fundamentos e instrumentos, preparar scripts de análisis y escribir código \LaTeX.\ Se reconoce el uso de generación sustantiva de contenido y de programación, no solamente corrección ortográfica. La comprobación automatizada comprende ejecución de pruebas de validación y compilación del documento; la evaluación académica de su pertinencia corresponde al equipo. No se atribuyen usos de la herramienta a los otros integrantes sin su confirmación.
}{}
\IfFileExists{secciones/08_consentimiento.tex}{\section{Protocolo de consentimiento informado}
\label{sec:consentimiento-kano}
\begingroup
\small
\setlength{\parskip}{3pt}

\noindent\textbf{Identificación institucional y académica.} El presente protocolo corresponde al equipo MRS, integrado por Mendoza Palma Allan Jeremy, Rochina Rochina Freddy David y Sánchez Centeno Roselyn Andreina, estudiantes de Ingeniería de Software de la Universidad Técnica Estatal de Quevedo. Se desarrolla en el marco de la asignatura Ingeniería de Requerimientos (ISR-401), bajo la docencia de Gleiston Cicerón Guerrero Ulloa, PhD, en la actividad FORM-SIM-01 del período SPA 2026--2027.

\noindent\textbf{Propósito.} Se solicita la participación en la aplicación de un cuestionario basado en el modelo de Kano para priorizar ocho características del dominio ``Unidad educativa: matrícula y calificaciones''. La actividad tiene fines académicos y formativos, no genera calificación para quien participa y sirve de preparación para la actividad individual TA-IND-02.

\noindent\textbf{Qué se solicita.} Cada persona responde, para cada una de las ocho características, dos preguntas: una sobre cómo se sentiría si la característica estuviera presente y otra sobre cómo se sentiría si estuviera ausente. Cada respuesta se elige entre cinco opciones. La duración estimada es de diez a quince minutos por persona; esta estimación queda sujeta a confirmación al definirse la modalidad de aplicación.

\noindent\textbf{Voluntariedad y retiro.} La participación es voluntaria. La persona puede omitir cualquier pregunta, interrumpir la sesión o retirarse sin consecuencias académicas ni personales. Si retira su consentimiento, el equipo deja de utilizar las respuestas correspondientes y gestiona su exclusión del análisis.

\noindent\textbf{Seudonimizado y tratamiento de la información.} En los archivos del repositorio público las personas participantes se identifican solo con los códigos P01 a P05. No se publican nombres, cédulas, teléfonos, correos electrónicos, audios, videos ni imágenes identificables. El uso de códigos constituye una seudonimización: reduce la exposición directa, pero no garantiza anonimato absoluto si existen datos adicionales que permitan establecer una vinculación. La información se emplea exclusivamente para el análisis académico de la actividad y se trata conforme a los principios aplicables de la normativa ecuatoriana de protección de datos personales~\cite{lopdp2021} y del Código de Ética y Conducta Profesional de la ACM~\cite{acm2018}. Este modelo no interpreta disposiciones legales concretas.

\noindent\textbf{Almacenamiento y acceso.} Los formularios firmados y cualquier evidencia que contenga datos personales se conservan únicamente en una carpeta institucional privada vinculada a cuentas institucionales de la UTEQ, con acceso restringido al equipo y al docente. En el repositorio público permanece solo el modelo sin completar, el archivo de respuestas seudonimizado y el manifiesto de evidencias. La conservación se limita al tiempo necesario para la evaluación académica. El presente modelo no implica aprobación institucional previa ni garantía jurídica absoluta.

\noindent\textbf{Manifestación del consentimiento.} Antes de iniciar, quien aplica el instrumento explica las condiciones anteriores y solicita la aceptación de forma verbal o escrita. La aceptación permite comenzar el cuestionario; la negativa impide su aplicación. Si se utiliza el formulario escrito, se completa en la copia privada y no se publica.

\noindent\textbf{Registro de la aplicación.} Para la aplicación del 27 de noviembre de 2026, Rochina comunicó que las personas participantes aceptaron el consentimiento. El repositorio dispone únicamente de las respuestas seudonimizadas exportadas: cinco formularios con los códigos P01, P04, P04, P02 y P05. No se cuenta con copias de los originales ni con registros que permitan confirmar si la aceptación fue verbal o escrita, la duración efectiva de cada sesión o los roles por código. Por tal motivo, el manifiesto público no incorpora filas ni huellas SHA-256 y no se publica información personal de quienes participaron.

\vspace{4pt}
\begin{center}
\renewcommand{\arraystretch}{1.8}
\begin{tabular}{@{}>{\bfseries\raggedright\arraybackslash}p{0.31\linewidth}>{\raggedright\arraybackslash}p{0.60\linewidth}@{}}
Nombre de quien participa & \rule[-2pt]{0.95\linewidth}{0.4pt} \\
Fecha & \rule[-2pt]{0.95\linewidth}{0.4pt} \\
Modalidad del consentimiento & Verbal\quad \fbox{\rule{0pt}{1.6ex}\rule{1.6ex}{0pt}}\quad Escrita\quad \fbox{\rule{0pt}{1.6ex}\rule{1.6ex}{0pt}} \\
Firma (solo si es escrita) & \rule[-2pt]{0.95\linewidth}{0.4pt} \\
Quien aplica el instrumento & \rule[-2pt]{0.95\linewidth}{0.4pt} \\
\end{tabular}
\end{center}
\endgroup
}{}

\nocite{pohl2025,wiegers2013,ireb2026,iso29148,lopdp2021,acm2018}
\bibliographystyle{IEEEtran}
\bibliography{referencias}
\end{document}
